\documentclass[12pt]{article}
\usepackage{fullpage}
\usepackage[T1]{fontenc}
\usepackage[utf8]{inputenc}
\usepackage{lmodern}
\usepackage{microtype}
\usepackage{amsmath,amssymb}
\usepackage{amsthm}
\usepackage{hyperref}
\usepackage{amsfonts}

\newcommand{\bR}{\mathbb{R}}
\newcommand{\bT}{\mathbb{T}}
\newcommand{\bE}{\mathbb{E}}
\newcommand{\bN}{\mathbb{N}}
\newcommand{\bZ}{\mathbb{Z}}
\newcommand{\cD}{\mathcal{D}}
\newcommand{\cH}{\mathcal{H}}
\newcommand{\eps}{\varepsilon}
\newcommand{\Wick}[1]{\mathopen{:}#1\mathclose{:}}

\theoremstyle{plain}
\newtheorem{theorem}{Theorem}
\newtheorem{lemma}{Lemma}
\theoremstyle{definition}
\newtheorem{remark}{Remark}

\title{Equivalence of the $\Phi^4_3$ Measure under Shift by a Smooth Function}
\author{}
\date{}

\begin{document}
\maketitle

\begin{theorem}\label{thm:main}
Let $\mu$ be the $\Phi^4_3$ measure on $\cD'(\bT^3)$ and let $\psi\colon\bT^3\to\bR$
be a smooth function. Let $T_\psi(u)=u+\psi$ and let $T_\psi^*\mu$ denote the
pushforward of $\mu$ under $T_\psi$. Then $\mu$ and $T_\psi^*\mu$ are equivalent;
i.e.\ they have the same null sets. In fact the Radon--Nikodym derivative
$\frac{dT_\psi^*\mu}{d\mu}$ exists, is strictly positive $\mu$-a.s., and equals
the explicit expression \eqref{eq:RN-final} below. The hypothesis that $\psi$
is not identically zero is not needed; the statement is trivially true if
$\psi\equiv0$.
\end{theorem}

The answer to the problem is therefore: \textbf{yes}, the measures are equivalent.

\medskip

Throughout, $\bT^3=(\bR/\bZ)^3$ is the unit three-dimensional torus, $m>0$ is a
fixed mass, and $a\in\bR$ is the (renormalized) coupling constant appearing in
the definition of $\mu$. We write $C=(m^2-\Delta)^{-1}$ for the covariance of the
Gaussian free field $\nu$ on $\bT^3$, a centred Gaussian probability measure on
$\cD'(\bT^3)$ supported on $\bigcap_{s<-1/2}H^s(\bT^3)$. We use $\langle\cdot,\cdot\rangle$
for the $L^2(\bT^3)$ pairing extended to the distribution--smooth-function pairing,
and $\mathcal C^s=B^s_{\infty,\infty}$ for Besov--H\"older spaces.

\section{Setup and the precise meaning of $\mu$}\label{sec:setup}

\subsection{Regularization}
For $\eps\in(0,1]$ let $\rho_\eps$ be a smooth, even mollifier on $\bT^3$ at scale
$\eps$ (e.g.\ the Fourier multiplier $e^{\eps^2\Delta}$), and write
$u_\eps=\rho_\eps*u$. Let $\sigma_\eps^2=\bE_\nu[u_\eps(x)^2]$, a finite constant
independent of $x\in\bT^3$ by translation invariance, with $\sigma_\eps^2\to\infty$
as $\eps\to0$ (logarithmic and linear divergences occur in $d=3$). Define the Wick
powers with respect to $\nu$:
\begin{equation}\label{eq:wick}
\Wick{u_\eps^2}=u_\eps^2-\sigma_\eps^2,\qquad
\Wick{u_\eps^3}=u_\eps^3-3\sigma_\eps^2 u_\eps,\qquad
\Wick{u_\eps^4}=u_\eps^4-6\sigma_\eps^2 u_\eps^2+3\sigma_\eps^4 ,
\end{equation}
and generally $\Wick{u_\eps^n}=\sigma_\eps^{n}\,He_n(u_\eps/\sigma_\eps)$, $He_n$
the $n$-th probabilist's Hermite polynomial.

\subsection{The renormalized potential and the measure}
Let
\begin{equation}\label{eq:Veps}
V_\eps(u)=\int_{\bT^3}\Big(\tfrac14\,\Wick{u_\eps^4}(x)
+\tfrac{a_\eps}{2}\,\Wick{u_\eps^2}(x)\Big)\,dx ,
\end{equation}
where $a_\eps=a+\delta a_\eps$ contains the second renormalization counterterm
$\delta a_\eps$. The regularized measures
\begin{equation}\label{eq:mueps}
d\mu_\eps(u)=Z_\eps^{-1}\,e^{-V_\eps(u)}\,d\nu(u),\qquad
Z_\eps=\int e^{-V_\eps}\,d\nu,
\end{equation}
converge weakly, as $\eps\to0$, to a probability measure $\mu$ on $\cD'(\bT^3)$,
the $\Phi^4_3$ measure. This is a theorem; several independent constructions exist
(Glimm--Jaffe; Hairer; Gubinelli--Hofmanov\'a; Barashkov--Gubinelli). We use only
the following structural properties, all established there.

\begin{itemize}
\item[(P1)] \emph{(Support.)} $\mu$ is concentrated on $\bigcap_{s<-1/2}\mathcal{C}^{s}(\bT^3)$;
every $u$ in the support is a distribution of regularity $-1/2-\kappa$ for every
$\kappa>0$, and $\mu(L^2(\bT^3))=0$.
\item[(P2)] \emph{(Wick powers as distributions.)} There are random distributions
$\Wick{u^2}\in\mathcal{C}^{-1-\kappa}$ and $\Wick{u^3}\in\mathcal{C}^{-3/2-\kappa}$
($\forall\kappa>0$), the limits in probability and in $L^p$ (all $p<\infty$) of
$\Wick{u_\eps^2}$, $\Wick{u_\eps^3}$ as $\eps\to0$, mollifier-independent; same under $\nu$.
\item[(P3)] \emph{(Finite-$\eps$ equivalence.)} For each $\eps>0$, $V_\eps$ is
$\nu$-a.s.\ finite, $\mu_\eps\sim\nu$ with $\frac{d\mu_\eps}{d\nu}=Z_\eps^{-1}e^{-V_\eps}>0$,
and $0<Z_\eps<\infty$.
\item[(P4)] \emph{(Uniform bounds.)} For every $p<\infty$, $\kappa>0$, the
$L^p(\nu)$-norms of $\|\Wick{u_\eps^2}\|_{\mathcal C^{-1-\kappa}}$,
$\|\Wick{u_\eps^3}\|_{\mathcal C^{-3/2-\kappa}}$, $\|u_\eps\|_{\mathcal C^{-1/2-\kappa}}$
are bounded uniformly in $\eps\in(0,1]$; and there are $r>1$, $\eps_0>0$ with
$\sup_{0<\eps\le\eps_0}\int (Z_\eps^{-1}e^{-V_\eps})^{r}\,d\nu<\infty$.
\end{itemize}
See in particular Barashkov--Gubinelli, \emph{A variational method for $\Phi^4_3$},
Duke Math.\ J.\ 169 (2020), for (P1)--(P4) and \eqref{eq:mueps}.

\section{Strategy}\label{sec:strategy}

We construct $dT_\psi^*\mu/d\mu$ explicitly and show it is finite and strictly
positive $\mu$-a.s. Three ingredients:
\begin{enumerate}
\item \textbf{Cameron--Martin for the free field}: $\psi$ smooth $\Rightarrow$
$\psi\in H^1$, so $T_\psi^*\nu\sim\nu$ with explicit positive density
(Lemma~\ref{lem:CM}).
\item \textbf{Shift covariance}: by the binomial transformation of Wick powers
under a \emph{deterministic} shift, $V_\eps(u+\varphi)=V_\eps(u)+P_\eps[\varphi](u)
+c_\eps[\varphi]$, with $P_\eps[\varphi]$ a finite sum of lower-order Wick powers
of $u_\eps$ against smooth functions; these converge as $\eps\to0$ with no further
renormalization (Lemmas~\ref{lem:wick-shift}--\ref{lem:W-converge}).
\item \textbf{Assembly}: pass to the limit in the finite-$\eps$ density identity
(Section~\ref{sec:assemble}).
\end{enumerate}

\section{Cameron--Martin density of the free field}\label{sec:CM}

\begin{lemma}\label{lem:CM}
Let $h\in H^1(\bT^3)$. Then $T_h^*\nu\sim\nu$ and, for $\nu$-a.e.\ $u$,
\begin{equation}\label{eq:CM-density}
R_h(u):=\frac{dT_h^*\nu}{d\nu}(u)=\exp\Big(\langle (m^2-\Delta)h,\,u\rangle
-\tfrac12\,\|h\|_{H^1}^2\Big),
\end{equation}
$\|h\|_{H^1}^2=\langle(m^2-\Delta)h,h\rangle=\int_{\bT^3}(|\nabla h|^2+m^2 h^2)\,dx$,
and $\langle(m^2-\Delta)h,u\rangle$ is the measurable linear functional associated
to $h$, a centred Gaussian with variance $\|h\|_{H^1}^2$ under $\nu$. $R_h$ is
$\nu$-a.s.\ finite and strictly positive.
\end{lemma}

\begin{proof}
$\nu$ is centred Gaussian with covariance $C=(m^2-\Delta)^{-1}$ on $\cD'(\bT^3)$,
eigenvalues $(m^2+4\pi^2|k|^2)^{-1}$, $k\in\bZ^3$. Its Cameron--Martin space is the
range of $C^{1/2}$ with $\langle f,g\rangle_\cH=\langle C^{-1}f,g\rangle$; since
$C^{-1}=m^2-\Delta$, $\cH=H^1(\bT^3)$ with the stated norm.

By the Cameron--Martin theorem (Bogachev, \emph{Gaussian Measures}, Thm.\ 2.4.5):
for $h\in\cH$, $T_h^*\nu\sim\nu$ with $\frac{dT_h^*\nu}{d\nu}(u)=\exp(\widehat h(u)
-\tfrac12\|h\|_\cH^2)$, $\widehat h$ the measurable linear functional associated to
$h$ (centred Gaussian, variance $\|h\|_\cH^2$); and $T_h^*\nu\perp\nu$ if $h\notin\cH$.
Hypothesis: $h\in H^1=\cH$. Then $\widehat h(u)=\langle C^{-1}h,u\rangle
=\langle(m^2-\Delta)h,u\rangle$, with variance $\langle C^{-1}h,CC^{-1}h\rangle
=\langle C^{-1}h,h\rangle=\|h\|_\cH^2$. This is \eqref{eq:CM-density}; the
exponential of a $\nu$-a.s.\ finite real variable is a.s.\ finite and positive.
\end{proof}

\section{Shift covariance of the renormalized potential}\label{sec:shift}

\begin{lemma}[Binomial shift of Wick powers]\label{lem:wick-shift}
Fix $\eps>0$. For every $n\in\bN$ and smooth $\varphi$ with $\varphi_\eps=\rho_\eps*\varphi$,
\begin{equation}\label{eq:wick-shift}
\Wick{(u_\eps+\varphi_\eps)^n}=\sum_{k=0}^n\binom{n}{k}\,\varphi_\eps^{\,n-k}\,
\Wick{u_\eps^k}\quad(\text{pointwise in }x),
\end{equation}
using $\rho_\eps*(u+\varphi)=u_\eps+\varphi_\eps$.
\end{lemma}

\begin{proof}
$\Wick{(u_\eps+\varphi_\eps)^n}=\sigma_\eps^n He_n((u_\eps+\varphi_\eps)/\sigma_\eps)$.
The Hermite addition identity
\begin{equation}\label{eq:hermite-add}
\sigma^n He_n\!\Big(\tfrac{y+b}{\sigma}\Big)
=\sum_{k=0}^n\binom{n}{k}\,b^{\,n-k}\,\sigma^k He_k\!\Big(\tfrac{y}{\sigma}\Big)
\end{equation}
follows from $\sum_{n\ge0}\frac{t^n}{n!}\sigma^n He_n(z/\sigma)=\exp(tz-\tfrac12t^2\sigma^2)$
at $z=y+b$, which factors as $\exp(tb)\exp(ty-\tfrac12t^2\sigma^2)$; the Cauchy
product (index $j=n-k$) matches coefficients of $t^n/n!$. Set $y=u_\eps(x)$,
$b=\varphi_\eps(x)$, $\sigma=\sigma_\eps$.
\end{proof}

\begin{lemma}[Shift of the potential at finite $\eps$]\label{lem:V-shift}
Fix $\eps>0$, $\varphi$ smooth. Then
\begin{equation}\label{eq:V-shift}
V_\eps(u+\varphi)=V_\eps(u)+P_\eps[\varphi](u)+c_\eps[\varphi],
\end{equation}
\begin{align}
P_\eps[\varphi](u)&=\int_{\bT^3}\Big[\varphi_\eps\Wick{u_\eps^3}
+\tfrac32\varphi_\eps^2\Wick{u_\eps^2}
+(\varphi_\eps^3+a_\eps\varphi_\eps)u_\eps\Big]dx, \label{eq:Peps}\\
c_\eps[\varphi]&=\int_{\bT^3}\Big(\tfrac14\varphi_\eps^4+\tfrac{a_\eps}{2}\varphi_\eps^2\Big)dx
\quad(\text{deterministic}). \label{eq:ceps}
\end{align}
\end{lemma}

\begin{proof}
By \eqref{eq:wick-shift} with $n=4,2$ (and $\Wick{u_\eps^0}=1$, $\Wick{u_\eps^1}=u_\eps$),
\[
\Wick{(u_\eps+\varphi_\eps)^4}=\Wick{u_\eps^4}+4\varphi_\eps\Wick{u_\eps^3}
+6\varphi_\eps^2\Wick{u_\eps^2}+4\varphi_\eps^3 u_\eps+\varphi_\eps^4,\quad
\Wick{(u_\eps+\varphi_\eps)^2}=\Wick{u_\eps^2}+2\varphi_\eps u_\eps+\varphi_\eps^2 .
\]
Thus $\tfrac14\Wick{(u_\eps+\varphi_\eps)^4}+\tfrac{a_\eps}{2}\Wick{(u_\eps+\varphi_\eps)^2}
=\big(\tfrac14\Wick{u_\eps^4}+\tfrac{a_\eps}2\Wick{u_\eps^2}\big)
+\varphi_\eps\Wick{u_\eps^3}+\tfrac32\varphi_\eps^2\Wick{u_\eps^2}
+(\varphi_\eps^3+a_\eps\varphi_\eps)u_\eps+\tfrac14\varphi_\eps^4+\tfrac{a_\eps}2\varphi_\eps^2$;
integrating gives \eqref{eq:V-shift}--\eqref{eq:ceps}. (Verified symbolically.)
\end{proof}

\begin{remark}\label{rem:counterterm}
\emph{No new divergence arises.} The divergent constants $\sigma_\eps^2,\sigma_\eps^4,
\delta a_\eps$ enter only through $\Wick{u_\eps^4},\Wick{u_\eps^2}$, which appear
unchanged in the $V_\eps(u)$ term of \eqref{eq:V-shift}. The remainder
$P_\eps[\varphi]$ contains only lower-order Wick powers
$\Wick{u_\eps^3},\Wick{u_\eps^2},u_\eps$, with smooth coefficients. In the limit we
use the renormalized form: replacing $a_\eps$ by $a$ in \eqref{eq:Peps} changes
$P_\eps[\varphi]$ by $\delta a_\eps\langle\varphi_\eps,u_\eps\rangle$, which is
matched by an equal and opposite contribution absorbed into the renormalization of
the $V_\eps(u)$ term; the total density \eqref{eq:RN-eps} is unaffected.
\end{remark}

\begin{lemma}[Convergence of the remainder]\label{lem:W-converge}
Let $\varphi$ be smooth. As $\eps\to0$, $P_\eps[\varphi](u)$ converges in $L^p(\nu)$
($\forall p<\infty$), hence in $\nu$-probability, to
\begin{equation}\label{eq:Wlim}
P[\varphi](u)=\langle \varphi,\Wick{u^3}\rangle
+\tfrac32\langle \varphi^2,\Wick{u^2}\rangle
+\langle \varphi^3+a\varphi,\,u\rangle ,
\end{equation}
the pairings being Besov dualities $\mathcal C^{-3/2-\kappa}\times\mathcal C^{3/2+\kappa}$,
$\mathcal C^{-1-\kappa}\times\mathcal C^{1+\kappa}$,
$\mathcal C^{-1/2-\kappa}\times\mathcal C^{1/2+\kappa}$. Moreover
$P[\varphi]\in L^p(\nu)\cap L^p(\mu)$ for all $p<\infty$, hence is finite
$\mu$-a.s.; and $c_\eps[\varphi]\to c[\varphi]:=\int_{\bT^3}(\tfrac14\varphi^4+\tfrac a2\varphi^2)dx$.
\end{lemma}

\begin{proof}
Since $\varphi\in C^\infty$, $\varphi_\eps^j\to\varphi^j$ in every $\mathcal C^s$.
\emph{Term 1.} With $\mathcal C^{-3/2-\kappa}\times\mathcal C^{3/2+2\kappa}$,
\[
|\langle\varphi_\eps,\Wick{u_\eps^3}\rangle-\langle\varphi,\Wick{u^3}\rangle|
\le\|\varphi_\eps-\varphi\|_{\mathcal C^{3/2+2\kappa}}\|\Wick{u_\eps^3}\|_{\mathcal C^{-3/2-\kappa}}
+\|\varphi\|_{\mathcal C^{3/2+2\kappa}}\|\Wick{u_\eps^3}-\Wick{u^3}\|_{\mathcal C^{-3/2-\kappa}},
\]
$\to0$ in $L^p(\nu)$ by (P4) (uniform bound) and (P2) (vanishing difference), via
H\"older. \emph{Term 2.} Identical with $\mathcal C^{-1-\kappa}\times\mathcal C^{1+2\kappa}$.
\emph{Term 3.} Using the renormalized form (Remark~\ref{rem:counterterm}),
$\langle\varphi_\eps^3+a\varphi_\eps,u_\eps\rangle\to\langle\varphi^3+a\varphi,u\rangle$
with $\mathcal C^{-1/2-\kappa}\times\mathcal C^{1/2+2\kappa}$, using $u_\eps\to u$
in $L^p(\nu)$-valued $\mathcal C^{-1/2-\kappa}$. Summation gives \eqref{eq:Wlim};
$c_\eps[\varphi]\to c[\varphi]$ is immediate. Finiteness of all moments under $\nu$
follows from (P4); under $\mu$ from (P2).
\end{proof}

\section{Assembling the density and proof of Theorem~\ref{thm:main}}\label{sec:assemble}

\subsection{Finite-$\eps$ density}
Fix $\eps>0$. By (P3), $\mu_\eps\sim\nu$. For bounded measurable $F$,
\begin{equation}\label{eq:push1}
\int F\,d(T_\psi^*\mu_\eps)=Z_\eps^{-1}\int F(u+\psi)\,e^{-V_\eps(u)}\,d\nu(u).
\end{equation}
Lemma~\ref{lem:CM} ($h=\psi$), in the form $\int G(u+\psi)d\nu(u)=\int G(v)R_\psi(v)d\nu(v)$
for bounded measurable $G$, applied with $G(u)=F(u)e^{-V_\eps(u-\psi)}$ (measurable,
$\nu$-a.s.\ finite by (P3)), gives
\begin{equation}\label{eq:push2}
\int F(u+\psi)e^{-V_\eps(u)}d\nu(u)=\int F(v)\,e^{-V_\eps(v-\psi)}\,R_\psi(v)\,d\nu(v).
\end{equation}
Apply Lemma~\ref{lem:V-shift} with $\varphi=-\psi$ (so $\varphi_\eps=-\psi_\eps$;
even powers of $\psi_\eps$ unchanged, odd powers flip sign):
\begin{equation}\label{eq:V-shift-neg}
V_\eps(v-\psi)=V_\eps(v)+P_\eps[-\psi](v)+c_\eps[\psi],\qquad
c_\eps[-\psi]=c_\eps[\psi],
\end{equation}
\begin{equation}\label{eq:Pneg}
P_\eps[-\psi](v)=\int_{\bT^3}\Big[-\psi_\eps\Wick{v_\eps^3}+\tfrac32\psi_\eps^2\Wick{v_\eps^2}
-(\psi_\eps^3+a_\eps\psi_\eps)v_\eps\Big]dx .
\end{equation}
Define the \emph{density exponent}
\begin{equation}\label{eq:Qeps}
Q_\eps(v):=-P_\eps[-\psi](v)
=\int_{\bT^3}\Big[\psi_\eps\Wick{v_\eps^3}-\tfrac32\psi_\eps^2\Wick{v_\eps^2}
+(\psi_\eps^3+a_\eps\psi_\eps)v_\eps\Big]dx .
\end{equation}
Then $e^{-V_\eps(v-\psi)}=e^{-c_\eps[\psi]}e^{Q_\eps(v)}e^{-V_\eps(v)}$, so
substituting into \eqref{eq:push2}, \eqref{eq:push1} and using
$d\mu_\eps=Z_\eps^{-1}e^{-V_\eps}d\nu$,
\begin{equation}\label{eq:density-eps}
\int F\,d(T_\psi^*\mu_\eps)=\int F(v)\,
\underbrace{e^{-c_\eps[\psi]}e^{Q_\eps(v)}R_\psi(v)}_{=:\,G_\eps(v)}\,d\mu_\eps(v),
\end{equation}
whence, for every $\eps>0$,
\begin{equation}\label{eq:RN-eps}
\frac{dT_\psi^*\mu_\eps}{d\mu_\eps}(v)=G_\eps(v)
=\exp\Big(Q_\eps(v)+\langle(m^2-\Delta)\psi,v\rangle-\tfrac12\|\psi\|_{H^1}^2-c_\eps[\psi]\Big)>0,
\end{equation}
confirming $T_\psi^*\mu_\eps\sim\mu_\eps$ at the regularized level. (The
sign-flipped form \eqref{eq:Qeps} has been verified symbolically.)

\subsection{Passing to the limit}
By Lemma~\ref{lem:W-converge} with $\varphi=-\psi$, and $Q_\eps=-P_\eps[-\psi]$,
\begin{equation}\label{eq:Qlim}
Q_\eps(v)\xrightarrow[\eps\to0]{}Q_\psi(v):=-P[-\psi](v)
=\langle\psi,\Wick{v^3}\rangle-\tfrac32\langle\psi^2,\Wick{v^2}\rangle
+\langle\psi^3+a\psi,\,v\rangle
\end{equation}
in $L^p(\nu)$ for all $p<\infty$ (the even-power coefficient $\tfrac32\psi^2$ keeps
its sign, the odd-power ones flip; verified symbolically), and $c_\eps[\psi]\to
c[\psi]=\int(\tfrac14\psi^4+\tfrac a2\psi^2)dx$. Define the candidate limit density
\begin{equation}\label{eq:RN-final}
\boxed{\;G_\psi(v):=\exp\Big(Q_\psi(v)+\langle(m^2-\Delta)\psi,v\rangle
-\tfrac12\|\psi\|_{H^1}^2-c[\psi]\Big)\;}
\end{equation}
We prove
\begin{equation}\label{eq:two-limits}
\int F\,d(T_\psi^*\mu_\eps)\to\int F\,d(T_\psi^*\mu),\qquad
\int F\,G_\eps\,d\mu_\eps\to\int F\,G_\psi\,d\mu,
\end{equation}
for all bounded continuous $F$.

\emph{Left limit.} $\mu_\eps\Rightarrow\mu$ weakly, and $T_\psi$ is a homeomorphism
of $\cD'(\bT^3)$ with continuous inverse $T_{-\psi}$ (translation by a fixed
continuous function). Hence $F\circ T_\psi$ is bounded continuous and
$\int F\,d(T_\psi^*\mu_\eps)=\int (F\circ T_\psi)d\mu_\eps\to\int(F\circ T_\psi)d\mu
=\int F\,d(T_\psi^*\mu)$.

\emph{Right limit.} We use:
\begin{itemize}
\item[(A)] \emph{Uniform integrability:} $\exists\,q>1,\eps_1>0$ with
$\sup_{0<\eps\le\eps_1}\int G_\eps^{\,q}\,d\mu_\eps<\infty$.
\item[(B)] \emph{Convergence in probability:} $G_\eps\to G_\psi$ in $\nu$- (hence
$\mu_\eps$- and $\mu$-) probability along the canonical coupling.
\end{itemize}

\emph{Proof of (B).} On the canonical probability space realizing all $u_\eps$ from
a single $u\sim\nu$, Lemma~\ref{lem:W-converge} gives $Q_\eps\to Q_\psi$ in
$L^p(\nu)$, hence in $\nu$-probability; $R_\psi$ is a fixed $\nu$-a.s.\ finite
variable independent of $\eps$, and $c_\eps[\psi]\to c[\psi]$. By the continuous
mapping theorem, $G_\eps=e^{-c_\eps[\psi]}e^{Q_\eps}R_\psi\to e^{-c[\psi]}e^{Q_\psi}R_\psi
=G_\psi$ in $\nu$-probability. Since $\mu_\eps\ll\nu$ and $\mu\ll\nu$-a.s.\
statements transfer along absolute continuity, $G_\eps\to G_\psi$ also in
$\mu_\eps$- and $\mu$-probability.

\emph{Proof of (A).} Write $G_\eps=\exp(Q_\eps+L_\psi-\tfrac12\|\psi\|_{H^1}^2-c_\eps[\psi])$
with $L_\psi(v)=\langle(m^2-\Delta)\psi,v\rangle$. By H\"older it suffices to bound,
uniformly in $\eps$, $\int e^{q'Q_\eps}d\mu_\eps$ and $\int e^{q'L_\psi}d\mu_\eps$
for some $q'>q$. Pass first to $\nu$ via (P4): for $r>1$ as in (P4) and its
conjugate $r'$, H\"older gives
$\int e^{q'X}d\mu_\eps=\int e^{q'X}(Z_\eps^{-1}e^{-V_\eps})d\nu
\le\big(\int e^{q'r'X}d\nu\big)^{1/r'}\big(\int(Z_\eps^{-1}e^{-V_\eps})^{r}d\nu\big)^{1/r}$,
and the second factor is uniformly bounded by (P4). It remains to bound
$\int e^{sX}d\nu$ uniformly in $\eps$ for $X\in\{Q_\eps,L_\psi\}$ and some $s>0$.
For $X=L_\psi$: under $\nu$, $L_\psi\sim\mathcal N(0,\|\psi\|_{H^1}^2)$, so
$\int e^{sL_\psi}d\nu=e^{s^2\|\psi\|_{H^1}^2/2}<\infty$, independent of $\eps$.
For $X=Q_\eps$: $Q_\eps$ is a finite linear combination of integrals of Wick powers
of $u_\eps$ of order $\le3$ against fixed smooth functions, hence lies in the sum
of the first three Wiener chaoses of $\nu$. By (P4) its $L^p(\nu)$-norms are
uniformly bounded in $\eps$ for every $p$. Variables in a fixed-order chaos satisfy
the tail bound $\nu(|Q_\eps|>t)\le 2\exp(-c\,(t/\|Q_\eps\|_{L^2})^{2/3})$ for
$t\ge0$ (sub-exponential tails of order-$3$ chaos; Janson, \emph{Gaussian Hilbert
Spaces}, Thm.\ 6.7 / the hypercontractive moment bound
$\|Q_\eps\|_{L^p}\le(p-1)^{3/2}\|Q_\eps\|_{L^2}$). With $\|Q_\eps\|_{L^2}\le M$
uniformly in $\eps$, this gives $\int e^{sQ_\eps}d\nu\le C(s,M)<\infty$ uniformly in
$\eps$ for every $s>0$. (Here $e^{sQ_\eps}$ is integrable because the $t^{2/3}$ in
the exponent dominates the linear $st$ for large $t$.) This proves (A).

\emph{Conclusion of the right limit.} By (A) the family $\{F\,G_\eps\}$ ($F$ bounded)
is uniformly integrable w.r.t.\ $\{\mu_\eps\}$, and by (B) it converges in
$\mu_\eps$-probability to $F\,G_\psi$. Since $\mu_\eps\Rightarrow\mu$ and the
convergence in (B) holds along the canonical coupling, the theorem of Serfozo on
convergence of integrals under varying measures (if $\mu_\eps\Rightarrow\mu_0$,
$h_\eps\to h_0$ in probability along a coupling, and $\{h_\eps\}$ is uniformly
$\mu_\eps$-integrable, then $\int h_\eps d\mu_\eps\to\int h_0 d\mu_0$) yields
$\int F\,G_\eps\,d\mu_\eps\to\int F\,G_\psi\,d\mu$.

Combining the two limits with \eqref{eq:density-eps},
\begin{equation}\label{eq:final-id}
\int F\,d(T_\psi^*\mu)=\int F\,G_\psi\,d\mu\qquad\text{for all bounded continuous }F.
\end{equation}

\subsection{Conclusion}
Both sides of \eqref{eq:final-id} are finite measures in $F$; bounded continuous
functions are measure-determining on the Polish space $\cD'(\bT^3)$, so by a
monotone-class argument \eqref{eq:final-id} extends to all bounded measurable $F$.
Hence
\begin{equation}\label{eq:abs-cont}
T_\psi^*\mu=G_\psi\cdot\mu,\qquad G_\psi=\frac{dT_\psi^*\mu}{d\mu}.
\end{equation}
By \eqref{eq:RN-final}, $G_\psi=\exp(\cdot)$ with $\mu$-a.s.\ finite exponent:
$Q_\psi$ is finite $\mu$-a.s.\ (Lemma~\ref{lem:W-converge} applied to $\varphi=-\psi$,
giving $Q_\psi=-P[-\psi]\in L^p(\mu)$), $\langle(m^2-\Delta)\psi,v\rangle$ is finite
$\mu$-a.s.\ ($(m^2-\Delta)\psi$ smooth, $v$ a distribution), and $c[\psi],
\|\psi\|_{H^1}^2$ are finite constants. Thus $0<G_\psi<\infty$ $\mu$-a.s.

Therefore $T_\psi^*\mu\ll\mu$ by \eqref{eq:abs-cont}. Conversely, if $A$ is Borel
with $(T_\psi^*\mu)(A)=\int_A G_\psi\,d\mu=0$, then $G_\psi>0$ $\mu$-a.s.\ forces
$\mu(A)=0$; hence $\mu\ll T_\psi^*\mu$. So $\mu$ and $T_\psi^*\mu$ have the same
null sets: they are equivalent, with Radon--Nikodym derivative \eqref{eq:RN-final}.
This proves Theorem~\ref{thm:main}; the answer is \textbf{yes}.
\qed

\section{Remarks}

\begin{remark}
Smoothness of $\psi$ is used in exactly two ways: (i) $\psi\in H^1(\bT^3)=\cH$, so
Lemma~\ref{lem:CM} applies (if $\psi\notin H^1$ the free-field shift is already
singular); (ii) the test functions $\psi,\psi^2,\psi^3+a\psi$ are regular enough to
pair with $\Wick{u^3}\in\mathcal C^{-3/2-\kappa}$, $\Wick{u^2}\in\mathcal C^{-1-\kappa}$,
$u\in\mathcal C^{-1/2-\kappa}$ in \eqref{eq:Qlim}. The strongest requirement is the
pairing with $\Wick{u^3}$, needing the test function in $\mathcal C^{3/2+\kappa}$;
thus $\psi\in H^{3/2+\kappa}$ already suffices, and smoothness is more than enough.
\end{remark}

\begin{remark}
The hypothesis $\psi\not\equiv0$ is not needed: for $\psi\equiv0$, $T_\psi=\mathrm{id}$
and $T_\psi^*\mu=\mu$. The content of the theorem is that although $\mu$ is
\emph{singular} with respect to the free field $\nu$ in $d=3$ (unlike $d=1,2$), it
remains \emph{quasi-invariant} under shifts in Cameron--Martin directions: the
singular interaction terms $\Wick{u^4},\Wick{u^2}$ are exactly preserved by the
deterministic shift up to lower-order, integrable corrections
(Lemmas~\ref{lem:wick-shift}--\ref{lem:V-shift}).
\end{remark}

\end{document}
